\documentclass[11pt,twocolumn]{article}

% ── Packages ──────────────────────────────────────────────
\usepackage[utf8]{inputenc}
\usepackage[T1]{fontenc}
\usepackage{amsmath,amssymb,amsfonts}
\usepackage{graphicx}
\usepackage{booktabs}
\usepackage{hyperref}
\usepackage{algorithm}
\usepackage{algpseudocode}
\usepackage{listings}
\usepackage{xcolor}
\usepackage[margin=0.75in]{geometry}
\usepackage{abstract}
\usepackage{titlesec}

% ── Code listing style ───────────────────────────────────
\definecolor{codebg}{RGB}{245,245,245}
\definecolor{codegreen}{RGB}{40,160,80}
\definecolor{codegray}{RGB}{128,128,128}
\definecolor{codepurple}{RGB}{160,32,240}

\lstdefinestyle{cppstyle}{
    backgroundcolor=\color{codebg},
    commentstyle=\color{codegreen},
    keywordstyle=\color{codepurple}\bfseries,
    numberstyle=\tiny\color{codegray},
    stringstyle=\color{red!60!black},
    basicstyle=\ttfamily\scriptsize,
    breaklines=true,
    captionpos=b,
    keepspaces=true,
    numbers=left,
    numbersep=5pt,
    showspaces=false,
    showstringspaces=false,
    showtabs=false,
    tabsize=4,
    language=C++,
    frame=single,
    rulecolor=\color{gray!30},
}

% ── Hyperref setup ────────────────────────────────────────
\hypersetup{
    colorlinks=true,
    linkcolor=blue!70!black,
    citecolor=blue!70!black,
    urlcolor=blue!70!black,
}

% ── Title ─────────────────────────────────────────────────
\title{%
    \textbf{NeuralOS: A Modular, Bit-Linear Compute Framework\\
    for Decentralized Intelligence}
}
\author{
    Abdul Hafizurrahman bin Azhar
}
\date{February 19, 2026}

% ══════════════════════════════════════════════════════════
\begin{document}
\maketitle

% ── Abstract ──────────────────────────────────────────────
\begin{onecolabstract}
The deployment of frontier-class Large Language Models (LLMs) with high reasoning capabilities is currently restricted by the ``Memory Wall''---the physical limitations of RAM capacity and bandwidth on consumer hardware. This paper introduces \textbf{NeuralOS}, a novel inference and training architecture that dismantles the monolithic LLM into a decentralized filesystem of independent 1.58-bit expert modules. By synthesizing zero-shot MLPMoE upcycling, BitNet~b1.58 quantization, and a novel Temporal Locality (``Sticky'') Router, we demonstrate the theoretical and practical viability of executing 100B+ parameter-class reasoning models on standard 16\,GB RAM CPU systems equipped with NVMe storage. NeuralOS shifts the paradigm from ``fitting the model in memory'' to ``paging intelligence on demand,'' effectively democratizing access to Super-Intelligence.
\end{onecolabstract}

\vspace{0.5em}
\noindent\textbf{Keywords:} Large Language Models, Mixture-of-Experts, BitNet, Model Compression, Edge Inference, NVMe Offloading, Decentralized AI

% ── 1. Introduction ──────────────────────────────────────
\section{Introduction}

The current trajectory of Artificial Intelligence faces a hardware divergence crisis. While State-of-the-Art (SOTA) models have scaled to hundreds of billions of parameters to achieve emergent reasoning and ``World Model'' capabilities, the average consumer hardware remains stagnant at 16\,GB to 32\,GB of unified system memory. This gap has centralized intelligence into data centers, creating privacy risks and accessibility barriers.

Standard compression techniques, such as 4-bit quantization (GGUF), have reached a plateau. A 70B parameter model in 4-bit precision still requires ${\approx}40$\,GB of RAM, far exceeding the capacity of consumer laptops. To bridge this gap, we must abandon the monolithic ``load-everything'' architecture.

We propose \textbf{NeuralOS}, a framework that treats the LLM not as a static binary but as a dynamic operating system. NeuralOS leverages the high sequential throughput of modern NVMe SSDs (3--7\,GB/s) to serve as a ``Cold Memory'' tier. However, naive offloading fails due to latency constraints. To solve this, NeuralOS introduces three architectural innovations:

\begin{enumerate}
    \item \textbf{Modular Bit-Linear Experts:} Decomposing dense models into thousands of granular, 1.58-bit experts.
    \item \textbf{Sticky Routing:} A routing algorithm that enforces temporal locality, minimizing disk I/O by reusing experts across long token sequences.
    \item \textbf{Speculative Prefetching:} Decoupling prediction from execution to mask I/O latency.
\end{enumerate}

% ── 2. Core Architecture ─────────────────────────────────
\section{Core Architecture: The Bit-Linear Expert}

The fundamental unit of computation in NeuralOS is the \textbf{Bit-Linear Expert}. Unlike traditional floating-point matrices, these experts are designed for extreme density and CPU-native execution.

\subsection{BitNet b1.58 Quantization}

Standard INT4 quantization is insufficient for our density goals. NeuralOS adopts the BitNet~b1.58 paradigm~\cite{ma2024era}, where weights are constrained to a ternary set $W \in \{-1, 0, +1\}$.

\begin{itemize}
    \item \textbf{Information Density.} This representation uses ${\approx}1.58$ bits per parameter ($\log_2 3$). A 100B parameter model compresses to roughly 22\,GB, theoretically fitting entirely within the virtual memory space of a 32\,GB SSD swap file.

    \item \textbf{Compute Efficiency.} BitNet replaces expensive Floating Point Multiply-Add (FMA) operations with integer Addition/Subtraction:
    \begin{equation}
        y = W \cdot x \approx \sum(x_{\text{active}}) - \sum(x_{\text{inactive}})
        \label{eq:bitnet-compute}
    \end{equation}
    This results in a $2\text{--}4\times$ speedup on CPU instruction sets (AVX-512/AMX) compared to FP16, crucial for compensating for the lower parallelism of CPUs.
\end{itemize}

\subsection{MLPMoE: Zero-Shot Upcycling}

We avoid the prohibitive cost of training MoEs from scratch by upcycling existing dense models (e.g., Llama-3-70B). We utilize MLPMoE~\cite{zhang2022moefication} (Multilayer Perceptron Mixture-of-Experts) decomposition via tensor slicing.

\begin{itemize}
    \item \textbf{Decomposition.} The massive Feed-Forward Network (FFN) matrices ($W_{\text{up}}$, $W_{\text{down}}$, $W_{\text{gate}}$) are sliced into $N$ smaller ``Expert'' matrices.

    \item \textbf{Mathematical Identity:}
    \begin{equation}
        \text{FFN}_{\text{dense}}(\mathbf{x}) = \sum_{i=1}^{N} \text{Expert}_i(\mathbf{x})
        \label{eq:ffn-decomposition}
    \end{equation}
    By routing tokens to only the top-$k$ experts (where $k \ll N$), we convert a dense operation into a sparse one without initial loss of knowledge.
\end{itemize}

% ── 3. Sticky Router ─────────────────────────────────────
\section{The ``Sticky'' Router: Solving the I/O Bottleneck}

The primary failure mode of disk-offloaded MoEs is \emph{thrashing}. Standard routers select different experts for every token, generating random I/O read requests that overwhelm the NVMe drive. NeuralOS introduces \textbf{Sticky Routing} to enforce Temporal Locality.

\subsection{Temporal Locality Regularization}

We modify the router to penalize expert switching. We introduce a regularization term $\mathcal{L}_{\text{sticky}}$ that minimizes the variance of the gating distribution $G(\mathbf{x})$ over a window of time $t$:

\begin{equation}
    \mathcal{L}_{\text{sticky}} = \lambda \sum_{t=1}^{T} \left\| G(\mathbf{x}_t) - G(\mathbf{x}_{t-1}) \right\|^2
    \label{eq:stickiness}
\end{equation}

\begin{itemize}
    \item \textbf{Effect.} This forces the model to ``commit'' to a set of experts (e.g., ``Coding Experts'') for an entire sequence (e.g., a Python function).
    \item \textbf{I/O Impact.} Instead of loading experts every token (${\sim}50$\,ms), the system loads experts once per sentence ($2000+$\,ms amortization window). This brings the required bandwidth down from GB/s to MB/s, well within the limits of consumer NVMe drives.
\end{itemize}

\subsection{Flash-MoE: Speculative Prefetching}

To hide the latency of unavoidable expert switches (e.g., changing topics), NeuralOS employs a ``Lookahead Predictor.''

\begin{itemize}
    \item \textbf{The Oracle.} A tiny, quantized dense model resides permanently in RAM. It processes the input stream $K$ tokens ahead of the main generator.
    \item \textbf{Async DMA.} If the Oracle predicts that the ``Medical Expert'' will be needed in 10 tokens, it triggers an asynchronous Direct Memory Access (DMA) request (via \texttt{io\_uring} on Linux) to fetch the expert from SSD to RAM before the computation requires it.
\end{itemize}

% ── 4. Neural Kernel ──────────────────────────────────────
\section{System Implementation: The Neural Kernel}

NeuralOS operates as a user-space kernel managing a three-tier memory hierarchy.

\subsection{Memory Hierarchy}

\begin{table}[h]
\centering
\caption{NeuralOS three-tier memory hierarchy.}
\label{tab:memory-hierarchy}
\begin{tabular}{@{}llll@{}}
\toprule
\textbf{Tier} & \textbf{Name} & \textbf{Contents} & \textbf{Location} \\
\midrule
L1 & Hot    & Attention Sinks, activations    & CPU Cache     \\
L2 & Warm   & Active expert set + KV Cache    & 16\,GB System RAM \\
L3 & Cold   & Full 1000+ expert library       & NVMe SSD      \\
\bottomrule
\end{tabular}
\end{table}

\subsection{Inference Loop Logic}

The core inference loop (Algorithm~\ref{alg:inference}) orchestrates prediction, prefetching, routing, and bit-linear computation in a pipelined fashion.

\begin{algorithm}[h]
\caption{NeuralOS Inference Step}\label{alg:inference}
\begin{algorithmic}[1]
\Procedure{Step}{context}
    \State $F \gets \textsc{Predict}(\text{ctx}, k{=}10)$ \Comment{Oracle lookahead}
    \State \textsc{AsyncPrefetch}$(F)$ \Comment{NVMe $\to$ RAM}
    \State $E \gets \textsc{Route}(\text{ctx})$ \Comment{Sticky router}
    \State $\mathbf{o} \gets \mathbf{0}$
    \For{$e \in E$}
        \State $\mathbf{w} \gets \textsc{GetExpert}(e)$ \Comment{Block until loaded}
        \State $\mathbf{o} \gets \mathbf{o} + \textsc{BitNetForward}(\mathbf{w}, \text{ctx})$
    \EndFor
    \State \Return $\mathbf{o}$
\EndProcedure
\end{algorithmic}
\end{algorithm}

\noindent A reference C++ implementation is provided below:

\begin{lstlisting}[style=cppstyle, caption={NeuralKernel inference loop.}]
class NeuralKernel {
    void step(Context ctx) {
        // 1. Oracle Prediction (Zero Latency)
        auto future_experts =
            predictor.predict(ctx, lookahead=10);

        // 2. Prefetching (Async I/O)
        vmm.prefetch(future_experts);

        // 3. Current Routing
        auto current_experts = router.route(ctx);

        // 4. Execution (BitNet Add/Sub)
        Tensor output = 0;
        for (auto exp_id : current_experts) {
            auto weights = vmm.get_expert(exp_id);
            output += bitnet_kernel.forward(
                weights, ctx);
        }
        return output;
    }
};
\end{lstlisting}

% ── 5. Training on the Edge ───────────────────────────────
\section{Training on the Edge}

NeuralOS also enables fine-tuning of these massive models on the same 16\,GB hardware by integrating BAdam~\cite{luo2024badam} and GaLore~\cite{zhao2024galore}.

\begin{itemize}
    \item \textbf{Block Coordinate Descent (BAdam).} Instead of updating all parameters simultaneously (which requires massive optimizer states), BAdam updates the model one ``expert block'' at a time. This reduces the peak memory requirement to the size of a single expert (${\approx}100$\,MB).

    \item \textbf{Gradient Low-Rank Projection (GaLore).} For the router and shared attention layers that must remain in memory, GaLore projects gradients into a low-rank subspace, reducing optimizer memory usage by up to 65\%.
\end{itemize}

% ── 6. Feasibility Analysis ──────────────────────────────
\section{Feasibility Analysis}

We present a feasibility comparison for running a 100B parameter model on a 16\,GB RAM laptop with a Gen4 NVMe SSD.

\begin{table}[h]
\centering
\caption{Dense baseline vs.\ NeuralOS on 16\,GB consumer hardware.}
\label{tab:feasibility}
\begin{tabular}{@{}lll@{}}
\toprule
\textbf{Metric} & \textbf{Dense (INT4)} & \textbf{NeuralOS} \\
\midrule
Total Model Size  & 55\,GB (OOM)       & ${\sim}$22\,GB (Paged) \\
Active RAM Usage  & 55\,GB             & ${\sim}$6--8\,GB       \\
I/O Requirement   & N/A (Crash)        & ${\sim}$150\,MB/s      \\
Compute Ops       & FP16/INT4 Mul-Add  & INT2 Add-Sub           \\
Est.\ Speed       & 0\,tok/s           & 8--12\,tok/s           \\
\bottomrule
\end{tabular}
\end{table}

\noindent\textbf{Note:} 8--12 tokens per second exceeds human reading speed, making this a viable platform for real-time chat and reasoning.

% ── 7. Conclusion ─────────────────────────────────────────
\section{Conclusion}

NeuralOS dismantles the hardware barrier to Artificial General Intelligence. By re-architecting the Large Language Model from a monolithic static file into a dynamic, modular filesystem of 1.58-bit experts, and by governing their execution with Sticky Routing and Predictive Prefetching, we enable high-reasoning, high-context AI on widely available consumer hardware. This work lays the foundation for a decentralized AI future, where ``Super-Intelligence'' is not a service rented from the cloud, but a software library running locally on the edge.

% ── References ────────────────────────────────────────────
\bibliographystyle{plain}
\begin{thebibliography}{5}

\bibitem{ma2024era}
S.~Ma, H.~Wang, L.~Ma, L.~Wang, W.~Wang, S.~Huang, L.~Dong, R.~Wang, J.~Xue, and F.~Wei.
\newblock The Era of 1-bit LLMs: All Large Language Models are in 1.58 Bits.
\newblock \emph{arXiv preprint arXiv:2402.17764}, 2024.

\bibitem{zhang2022moefication}
Z.~Zhang, Y.~Lin, Z.~Liu, P.~Li, M.~Sun, and J.~Zhou.
\newblock MoEfication: Transformer Feed-forward Layers are Mixtures of Experts.
\newblock In \emph{Findings of ACL}, 2022.

\bibitem{zhao2024galore}
J.~Zhao, Z.~Zhang, B.~Chen, Z.~Wang, A.~Anandkumar, and Y.~Tian.
\newblock GaLore: Memory-Efficient LLM Training by Gradient Low-Rank Projection.
\newblock \emph{arXiv preprint arXiv:2403.03507}, 2024.

\bibitem{luo2024badam}
Q.~Luo, Y.~Xu, Z.~Liu, G.~Bian, J.~Cheng, and L.~Zhang.
\newblock BAdam: A Memory Efficient Full Parameter Training Method for Large Language Models.
\newblock \emph{arXiv preprint arXiv:2404.02827}, 2024.

\bibitem{wang2024bitnet}
H.~Wang, S.~Ma, L.~Dong, S.~Huang, H.~Wang, L.~Ma, F.~Yang, R.~Wang, Y.~Wu, and F.~Wei.
\newblock BitNet: Scaling 1-bit Transformers for Large Language Models.
\newblock \emph{arXiv preprint arXiv:2310.11453}, 2023.

\end{thebibliography}

\end{document}
